% ============================================================
% RÉSUMÉ
% ============================================================
\cleardoublepage
\chapter*{RÉSUMÉ}
\addcontentsline{toc}{chapter}{Résumé}

L'approvisionnement en eau potable constitue un enjeu majeur
pour les populations urbaines, particulièrement dans les zones
où les réseaux de distribution sont confrontés à des problèmes
de pertes d'eau, de variations de pression, de défaillances des
équipements et de difficultés de surveillance. Dans la commune
de Kadutu, à Bukavu, ces contraintes rendent nécessaire la mise
en place de solutions technologiques capables d'améliorer la
gestion et le suivi du réseau.

Le présent travail porte sur le \textbf{développement d'un
	jumeau digital intégrant l'intelligence artificielle pour
	l'optimisation du réseau de fourniture d'eau potable dans la
	commune de Kadutu}. L'objectif est de concevoir une représentation
numérique du réseau permettant de suivre son fonctionnement,
de simuler différents scénarios et d'identifier plus rapidement
certaines anomalies susceptibles d'affecter la distribution
de l'eau.

La solution proposée s'appuie sur une architecture informatique
intégrant une base de données, une interface de visualisation
géographique, une API développée avec \textbf{FastAPI}, ainsi
que des techniques d'intelligence artificielle destinées à
l'analyse des données du réseau. Le système permet notamment
de représenter les principaux éléments du réseau, de surveiller
certains paramètres, de simuler des situations de dysfonctionnement
et de faciliter la prise de décision.

Les résultats attendus montrent qu'un jumeau digital associé
à l'intelligence artificielle peut constituer un outil d'aide
à la gestion, à la surveillance et à la maintenance du réseau
d'eau potable. Une telle approche contribue à améliorer la
connaissance du réseau, à anticiper certaines anomalies et à
soutenir une gestion plus efficace des ressources en eau.

\vspace{0.5cm}

\noindent\textbf{Mots-clés :} Jumeau digital, intelligence
artificielle, réseau d'eau potable, optimisation, détection
d'anomalies, simulation, FastAPI, Kadutu.

\cleardoublepage